\section{환경 생성과 능동 실패 탐색}

\subsection*{장면 구조화와 후보 생성}
환경은 객체(목표물, 장애물, 목적지, 사람)의 위치·크기·회전·물성으로 이루어진 파라미터 벡터로 정의된다. 
로봇팔 환경은 그릇·접시 위치와 장애물 상자의 위치·크기·높이·회전 등 \textbf{10개 \gl{var}{변수}}, 
휴머노이드 환경은 상자·바닥 물성, 통로 폭, 장애물 2개의 형상, 사람의 이동을 포함한 \textbf{21개 변수}를 제공한다(표~\ref{tab:axes}). 
환경 생성기는 하나의 장면 파라미터에서 3가지 표현을 함께 만든다.
\begin{itemize}
\item \textbf{3D 장면 그래프}: 객체를 노드로, 거리·인접·목표 영역 점유 같은 공간 관계를 간선으로 기록한다.
\item \textbf{이동 지도}: 바닥을 격자로 나누고 장애물·상자가 차지하는 칸을 이동 불가로 표시한 점유 격자 지도로, 휴머노이드의 경로 계획에 쓰인다.
\item \textbf{물리 시뮬레이션 모델}: 객체의 형상·위치·질량·마찰을 MuJoCo의 MJCF~\cite{mjcf}로 기술한다.
\end{itemize}
예를 들어 휴머노이드 장면에서 상자를 통로 중앙으로 옮기면 MJCF의 상자 좌표, 이동 지도의 목표 부근 격자,
장면 그래프의 상자–목표 관계(``비어 있음''$\to$``점유'')가 함께 바뀐다. 
3개 표현들은 같은 장면 파라미터에서 만들어지므로 생성기는 실행 전에 서로 맞는지(예: 장면 그래프와 MJCF의 상자 위치) 확인하고,
물체끼리 겹치거나 작업 영역을 벗어나거나 로봇 몸과 겹치는 장면은 실행하지 않는다. 3개 표현들의 값은 시험 환경의 정답으로, 실제 로봇이라면 SLAM(Simultaneous Localization and Mapping)~\cite{slam}과 물체 검출로 얻어야 한다(4절). 

\subsection*{행동 피드백 기반 탐색 절차}
\begin{enumerate}
\item 이전 실행의 작업 성공 여부, 실제 행동, 로봇 기능(경로 계획, 상자 밀기 등)의 실행 결과, 접촉·이동량 측정, 종료 상태를 요약한다.
\item AFS의 LLM이 관측 사실을 근거로 실패 가설을 세우고, 허용 범위 안의 환경 변경값을 정해진 형식으로 제안한다.
\item 후보 검증기가 값의 범위, 근거로 든 이전 실행이 실제로 있는지, 바꾼 변수의 수, 남은 시도 횟수를 검사한다. 범위 밖 값은 실행 전에 거부한다.
\item 새 환경에서 로봇을 실행하고 독립된 평가기가 판정한다. 예상과 실제의 차이를 다음 제안에 반영한다.
\end{enumerate}
성공 기록도 함께 제공하여, 이미 알려진 실패만 반복하지 않고 성공 경계 근처를 탐색하도록 한다.
Random 테스트는 같은 범위에서 바꿀 변수와 그 값을 균등하게 무작위 추출하며 결과를 보고 재추출하지 않는다.

\begin{table}[!t]
\centering\footnotesize\setlength{\tabcolsep}{3pt}
\caption{AFS의 입력 관측과 출력 형식}\label{tab:afsio}
\begin{tabularx}{\linewidth}{@{}p{24mm}Y@{}}
\toprule
\multicolumn{2}{@{}l}{\textbf{입력: 이전 실행의 행동 관측}} \\
\midrule
성공 여부·종료 상태 & PASS/FAIL, 종료 사유, 목표 거리 변화 \\
로봇 기능 실행 결과 & 경로 계획 실패·반복, 중간 목표 도착, 밀기 등 조작 요청과 결과 \\
물리 측정값 & 접촉 시간·접촉력, 물체 이동량(로봇팔) \\
장면 관계 & 장면 그래프의 객체 간 거리, 목표 영역 점유 여부(예: 상자–목표 ``비어 있음'') \\
예상과 측정 & 이전 제안의 예상 유형과 실제 검출 유형 \\
\midrule
\multicolumn{2}{@{}l}{\textbf{출력: 제안 한 건의 구성}} \\
\midrule
근거 실행 & 기준으로 삼은 이전 실행의 식별자와 선택 이유 \\
변경 & 변경할 변수와 값(1--2개) \\
가설·취약점 & 관측에서 추론한 실패 원인과 대안 설명 \\
예상 실패 유형 & 12종 중 목표로 하는 유형 \\
반증 조건·측정 계획 & 가설이 틀렸다고 볼 결과와 확인할 측정값 \\
\bottomrule
\end{tabularx}
\end{table}

표~\ref{tab:afsio}는 AFS가 받는 관측과 내놓는 제안의 형식이다. 
제안마다 근거 실행과 반증 조건을 명시하게 하여, 후보 검증기가 근거로 든 실행이 실제로 있는지 확인하고 결과가 나온 뒤 예상과 실제를 비교할 수 있게 하였다. 
예상 유형은 판정에 쓰지 않으며, 실제 유형은 평가기의 측정으로만 정한다.
휴머노이드 본 실험에서 AFS 제안 1회는 중앙값 약 1분(34--93초)으로, 로봇 실행 1회(중앙값 약 2분)보다 짧았다.

\subsection*{실제 탐색 연결 예}
표~\ref{tab:link}는 두 실험에서 이전 실행의 관측이 LLM의 가설과 다음 환경 제안을 거쳐 실제 실패로 이어진 기록이다.
\begin{table*}[!t]
\centering\small
\caption{관측에서 다음 시험으로 이어진 실제 기록}\label{tab:link}
\begin{tabularx}{\textwidth}{@{}p{22mm}YY@{}}
\toprule
단계 & 로봇팔 (첫 번째 세트) & 휴머노이드 (예비 실험) \\
\midrule
이전 실행 & 시작 장면 FAIL: 로봇–장애물 접촉 13.3초, 그릇 미이동 & 시작 장면 PASS: 연속 이동으로 목표 도달 \\
LLM 가설 & 이송 경로 중앙의 더 크고 높은 상자가 접근을 막을 것 & 상자를 목표 쪽으로 옮기면 목표 영역이 점유될 것 \\
제안 & 장애물 폭·깊이 6\,cm, 높이 17\,cm & 상자를 통로 가장자리에서 목표 위치로 이동 (변수 1개 변경) \\
결과 & FAIL, 접촉 11.0초, 그릇 미이동 & FAIL, 목표 위치 점유 검출 \\
\bottomrule
\end{tabularx}
\end{table*}
